\textbf{Phần 1. Trắc nghiệm lựa chọn}
\setcounterpageref{ex}{0}

\begin{ex} %[4D8N42]
Tính chất nhân một số với một tổng được viết dưới dạng công thức là
\choice
{\True $a \times (b + c) = a \times b + a \times c$}
{$a \times (b + c) = a \times b + c$}
{$a \times (b + c) = a + b \times c$}
{$a \times (b + c) = a \times b \times c$}
\loigiai{
$a \times (b + c) = a \times b + a \times c$.\\
\textbf{Đáp án của câu hỏi là A}
}
\end{ex}

\begin{ex} %[4D8N42]
Tính chất nhân một số với một hiệu được viết dưới dạng công thức là
\choice
{$a \times (b - c) = a \times b - c$}
{$a \times (b - c) = a - b \times c$}
{$a \times (b - c) = a \times b : c$}
{\True $a \times (b - c) = a \times b - a \times c$}
\loigiai{
$a \times (b - c) = a \times b - a \times c$.\\
\textbf{Đáp án của câu hỏi là D}
}
\end{ex}

\begin{ex} %[4D8H42]
Tính thuận tiện: $35 \times 18 + 35 \times 82$.
\choice
{$3400$}
{\True $3500$}
{$3600$}
{$3700$}
\loigiai{
$35 \times (18 + 82) = 35 \times 100 = 3500$.\\
\textbf{Đáp án của câu hỏi là B}
}
\end{ex}

\begin{ex} %[4D8H42]
Tính thuận tiện: $48 \times 115 - 48 \times 15$.
\choice
{$4600$}
{$5000$}
{\True $4800$}
{$4700$}
\loigiai{
$48 \times (115 - 15) = 48 \times 100 = 4800$.\\
\textbf{Đáp án của câu hỏi là C}
}
\end{ex}

\begin{ex} %[4D8H42]
Biểu thức $125 \times 99 + 125$ có giá trị bằng
\choice
{\True $12\,500$}
{$12\,400$}
{$12\,600$}
{$13\,000$}
\loigiai{
$125 \times (99 + 1) = 125 \times 100 = 12\,500$.\\
\textbf{Đáp án của câu hỏi là A}
}
\end{ex}

\begin{ex} %[4D8H42]
Tính nhanh: $27 \times 11$.
\choice
{$287$}
{$307$}
{\True $297$}
{$277$}
\loigiai{
$27 \times (10 + 1) = 270 + 27 = 297$.\\
\textbf{Đáp án của câu hỏi là C}
}
\end{ex}

\begin{ex} %[4D8V42]
Tính nhanh giá trị biểu thức: $135 \times 45 + 135 \times 55$.
\choice
{$13\,400$}
{$13\,600$}
{$14\,000$}
{\True $13\,500$}
\loigiai{
$135 \times (45 + 55) = 135 \times 100 = 13\,500$.\\
\textbf{Đáp án của câu hỏi là D}
}
\end{ex}

\begin{ex} %[4D8V42]
Một cửa hàng ngày thứ nhất bán được 25 bao gạo, ngày thứ hai bán được 35 bao gạo. Mỗi bao gạo nặng $50\text{ kg}$. Tính tổng số ki-lô-gam gạo bán trong 2 ngày bằng cách thuận tiện nhất.
\choice
{$2800\text{ kg}$}
{\True $3000\text{ kg}$}
{$3200\text{ kg}$}
{$3500\text{ kg}$}
\loigiai{
$(25 + 35) \times 50 = 60 \times 50 = 3000\text{ kg}$.\\
\textbf{Đáp án của câu hỏi là B}
}
\end{ex}

\begin{ex} %[4D8V42]
Tính nhanh: $64 \times 101 - 64$.
\choice
{$6300$}
{$6500$}
{\True $6400$}
{$6464$}
\loigiai{
$64 \times (101 - 1) = 64 \times 100 = 6400$.\\
\textbf{Đáp án của câu hỏi là C}
}
\end{ex}

\begin{ex} %[4D8V42]
Một mảnh đất hình chữ nhật có chiều dài $28\text{ m}$, chiều rộng $12\text{ m}$. Nếu tăng chiều dài thêm $5\text{ m}$ thì diện tích mảnh đất tăng thêm bao nhiêu mét vuông?
\choice
{$50\text{ m}^2$}
{$70\text{ m}^2$}
{$80\text{ m}^2$}
{\True $60\text{ m}^2$}
\loigiai{
Diện tích tăng: $5 \times 12 = 60\text{ m}^2$.\\
\textbf{Đáp án của câu hỏi là D}
}
\end{ex}

\begin{ex} %[4D8V42]
Tính giá trị biểu thức: $45 \times 12 + 45 \times 18 + 45 \times 70$.
\choice
{\True $4500$}
{$4400$}
{$4600$}
{$5000$}
\loigiai{
$45 \times (12 + 18 + 70) = 45 \times 100 = 4500$.\\
\textbf{Đáp án của câu hỏi là A}
}
\end{ex}

\begin{ex} %[4D8C42]
Tính nhanh giá trị biểu thức: $A = 99 \times 45 + 45$.
\choice
{$4400$}
{\True $4500$}
{$4600$}
{$4455$}
\loigiai{
$A = 45 \times (99 + 1) = 45 \times 100 = 4500$.\\
\textbf{Đáp án của câu hỏi là B}
}
\end{ex}

\textbf{Phần 2. Trắc nghiệm đúng sai}
\setcounterpageref{ex}{0}

\begin{ex} %[4D8H42]
Xét tính đúng sai của các công thức phân phối:
\choiceTF
{\True $a \times (b + c) = a \times b + a \times c$.}
{\True $a \times (b - c) = a \times b - a \times c$.}
{$(a + b) \times c = a \times c + b$.}
{\True $(a + b) \times c = a \times c + b \times c$.}
\loigiai{
\begin{itemize}
    \item $a \times (b + c) = a \times b + a \times c$. Đúng.
    \item $a \times (b - c) = a \times b - a \times c$. Đúng.
    \item Phải nhân cả $b$ với $c$. Ý 3 Sai, Ý 4 Đúng.
\end{itemize}
\begin{itemize}
    \item a) \textbf{Đúng}.
    \item b) \textbf{Đúng}.
    \item c) \textbf{Sai}.
    \item d) \textbf{Đúng}.
\end{itemize}
}
\end{ex}

\begin{ex} %[4D8V42]
Hai tổ may khẩu trang: Tổ 1 may được 140 hộp, tổ 2 may được 160 hộp. Mỗi hộp chứa 50 chiếc khẩu trang. Xét tính đúng sai:
\choiceTF
{\True Cả hai tổ may được tất cả 300 hộp khẩu trang.}
{\True Tổ 2 may nhiều hơn tổ 1 là 1000 chiếc khẩu trang.}
{Tổng số khẩu trang cả hai tổ may được là 12000 chiếc.}
{\True Tổng số khẩu trang cả hai tổ may được là 15000 chiếc.}
\loigiai{
\begin{itemize}
    \item $140 + 160 = 300$ hộp: Đúng.
    \item $(160 - 140) \times 50 = 20 \times 50 = 1000$ chiếc: Đúng.
    \item Tổng: $300 \times 50 = 15\,000$ chiếc $\neq 12\,000$ chiếc. Ý 3 Sai, Ý 4 Đúng.
\end{itemize}
\begin{itemize}
    \item a) \textbf{Đúng}.
    \item b) \textbf{Đúng}.
    \item c) \textbf{Sai}.
    \item d) \textbf{Đúng}.
\end{itemize}
}
\end{ex}

\textbf{Phần 3. Trả lời ngắn}
\setcounterpageref{ex}{0}

\begin{ex} %[4D8H42]
Tính thuận tiện: $45 \times 38 + 45 \times 62$.
\shortans{4500}
\loigiai{
\textbf{Đáp số:} $4500$\\
$45 \times 100 = 4500$.
}
\end{ex}

\begin{ex} %[4D8H42]
Tính thuận tiện: $78 \times 125 - 78 \times 25$.
\shortans{7800}
\loigiai{
\textbf{Đáp số:} $7800$\\
$78 \times 100 = 7800$.
}
\end{ex}

\begin{ex} %[4D8V42]
Tính nhanh: $56 \times 99 + 56$.
\shortans{5600}
\loigiai{
\textbf{Đáp số:} $5600$\\
$56 \times 100 = 5600$.
}
\end{ex}

\begin{ex} %[4D8V42]
Tính nhanh: $84 \times 101 - 84$.
\shortans{8400}
\loigiai{
\textbf{Đáp số:} $8400$\\
$84 \times 100 = 8400$.
}
\end{ex}

\begin{ex} %[4D8V42]
Một lớp học mua 15 quyển vở Toán và 15 quyển vở Tiếng Việt cùng giá $8000$ đồng một quyển. Tính tổng số tiền mua vở theo nghìn đồng (chỉ ghi số nghìn đồng).
\shortans{240}
\loigiai{
\textbf{Đáp số:} $240$\\
$15 \times (8 + 8) = 240$ nghìn đồng.
}
\end{ex}

\begin{ex} %[4D8C42]
Tính nhanh: $12 \times 35 + 12 \times 65 + 24 \times 50$.
\shortans{2400}
\loigiai{
\textbf{Đáp số:} $2400$\\
$12 \times (35 + 65) + 12 \times 100 = 12 \times 100 + 12 \times 100 = 2400$.
}
\end{ex}
